\subsection{Completing the Square} 
\begin{fact}[Claim: We can ``complete'' a square]
For any expression $x^2+bx$, we can find a constant $c$ such that $x^2+bx+c$ is a perfect square trinomial.
\end{fact}
Thinking of multiplication as representing the area of a rectangle, we can think of the constant $c$ as the ``missing piece'' of a square.
\begin{example}
Find the constant $c$ such that $x^2+6x+c$ is a perfect square, and write it as such.
\begin{lpic}{}{4}
\regtextleft{0.25}{-1}{$x^2+6x+$};
\regtextleft{4}{-1}{${}=(x$};
\regtext{6}{-1}{$+$};
\regtextleft{7}{-1}{$)^2$};
\draw (2*\value{cols}-3,-1) rectangle (2*\value{cols},-4);
\draw (2*\value{cols}-3,-2.75) -- ++(3,0);
\draw (2*\value{cols}-1.25,-1) -- ++(0,-3);
\end{lpic}
\end{example}
\noi Sometimes, fractions appear in the constant.
\begin{example}
Complete the square.
\begin{lpic}{}{4}
\regtextleft{0.25}{-1}{$x^2-5x+$};
\regtextleft{4}{-1}{${}=(x$};
\regtext{6}{-1}{$+$};
\regtextleft{7}{-1}{$)^2$};
\draw (2*\value{cols}-3,-1) rectangle (2*\value{cols},-4);
\draw (2*\value{cols}-3,-2.75) -- ++(3,0);
\draw (2*\value{cols}-1.25,-1) -- ++(0,-3);
\end{lpic}
\end{example}
\noi We can do this even when the number $b$ is already a fraction.
\begin{example}
Complete the square.
\begin{lpic}{}{4}
\regtextleft{0.25}{-1}{$x^2+\frac{5}{3}x+$};
\regtextleft{4}{-1}{${}=(x$};
\regtext{6}{-1}{$+$};
\regtextleft{7}{-1}{$)^2$};
\draw (2*\value{cols}-3,-1) rectangle (2*\value{cols},-4);
\draw (2*\value{cols}-3,-2.75) -- ++(3,0);
\draw (2*\value{cols}-1.25,-1) -- ++(0,-3);
\end{lpic}
\end{example}
